\section{Workspace and Communication Windows}
Resident state explains what is already sitting on the rank. Workspace explains
what spikes only while a phase is running.

\begin{quote}
The estimator treats workspace as a \emph{live window}, not as a sum over every
temporary allocation that might appear somewhere in the phase.
\end{quote}

That choice is deliberate. Most kernels do not keep every temporary alive at
the same time, so summing all of them would be too pessimistic.

\subsection{Forward and Backward Workspace}
There are four main sources of phase-local workspace:
\begin{itemize}[leftmargin=1.5em]
\item attention workspace,
\item MLP or backward-kernel workspace,
\item communication workspace from tensor or sequence parallelism,
\item loss materialization workspace.
\end{itemize}

\paragraph{Attention backend choice.}
The estimator uses two different attention workspace models:
\begin{itemize}[leftmargin=1.5em]
\item \textbf{standard attention} scales with local tokens, local heads, and
effective key length because that path really does materialize a larger
score-like object,
\item \textbf{sdpa} and \textbf{flash2} use a width-based proxy tied to the
largest local attention-path width, because those kernels avoid the eager
score-matrix behavior the standard path exhibits.
\end{itemize}

This is an important modeling choice: the estimator is not saying all attention
backends behave the same and then patching them afterward. It starts from the
fact that the kernel strategies are different.

The two compact formulas are:
\begin{align*}
W_{\mathrm{attn,fwd}}^{\mathrm{standard}}
&\propto
T_{\mathrm{local}} \cdot h_{\mathrm{local}} \cdot S_{\mathrm{eff}}, \\
W_{\mathrm{attn,fwd}}^{\mathrm{non\text{-}eager}}
&\propto
T_{\mathrm{local}} \cdot d_{\mathrm{attn,path}}.
\end{align*}
The exact constants differ, but the structural point is the important one:
standard attention scales with a score-like object, while \texttt{sdpa} and
\texttt{flash2} scale with the widest local attention path.

\paragraph{Checkpointing.}
Checkpointing shifts pressure from retained activations into backward-time
recompute workspace. In practice that means:
\begin{itemize}[leftmargin=1.5em]
\item forward usually gets cheaper,
\item backward gets an extra recompute term,
\item total peak may move phases rather than simply shrinking everywhere.
\end{itemize}

\paragraph{Parallel communication.}
If tensor parallel or sequence parallel is enabled, the estimator adds explicit
communication windows for those exchanges. These are local windows. They are
not global cluster totals.

\subsection{Optimizer Update Scratch}
The optimizer also needs its own temporary workspace. The simplest way to think
about it is:

\begin{quote}
How much temporary storage does one rank need while it is actively applying one
update to the trainable set it owns?
\end{quote}

For single-GPU and DDP, that scratch scales with the full local trainable set.
For sharded ZeRO paths, it may scale with the local shard instead.

The compact form is:
\[
W_{\mathrm{update}} \approx
\text{local update elements} \times \text{update bytes per element}.
\]

\subsection{Configured Bucket Objects}
Distributed training introduces another kind of workspace: reducer buckets,
all-gather buckets, and prefetch buckets. The estimator now models those as
explicit configured objects rather than hidden byte caps.

\paragraph{DDP.}
DDP has one reducer bucket. Its capacity is:
\[
B_{\mathrm{ddp}}^{\mathrm{cap}} =
E_{\mathrm{ddp}} \cdot b_g^\star,
\]
where $E_{\mathrm{ddp}}$ is the configured bucket size in gradient elements.

\paragraph{ZeRO.}
ZeRO has three relevant bucket objects:
\begin{itemize}[leftmargin=1.5em]
\item an all-gather bucket,
\item a reduce bucket,
\item a prefetch bucket.
\end{itemize}

Their capacities come from configured element counts as well. Because the
estimator is per-rank, it converts those distributed bucket sizes into local
rank capacities by dividing over the data-parallel group before multiplying by
the appropriate dtype width.

The capacity equations are:
\begin{align*}
B_{\mathrm{zero,ag}}^{\mathrm{cap}} &=
\dfrac{E_{\mathrm{zero}}}{D_{\mathrm{P}}} \cdot b_p^\star, \\
B_{\mathrm{zero,red}}^{\mathrm{cap}} &=
\dfrac{E_{\mathrm{zero}}}{D_{\mathrm{P}}} \cdot b_g^\star, \\
B_{\mathrm{zero,pf}}^{\mathrm{cap}} &=
\dfrac{E_{\mathrm{prefetch}}}{D_{\mathrm{P}}} \cdot b_p^\star.
\end{align*}
Again, the implementation rounds to whole elements, but the simplified form is
easier to read and captures the real regime difference: these capacities shrink
with the data-parallel degree.

\paragraph{Capacity versus live bytes.}
One subtle point matters here: the configured bucket size describes the
capacity of the object, not the amount of useful payload that always fills it.
The estimator therefore clips each bucket to the local payload that can
actually occupy it. A rank cannot actively hold more gradient bytes than it
owns, even if the configured reducer bucket is larger than that.

That idea is captured by:
\[
B_{\mathrm{active}} = \min(B_{\mathrm{capacity}}, M_{\mathrm{local\ payload}}).
\]

\subsection{DDP and ZeRO Windows In Words}
Once those bucket objects exist, the distributed windows become easier to read.

\paragraph{DDP.}
DDP contributes:
\begin{itemize}[leftmargin=1.5em]
\item the reducer bucket itself,
\item and a communication overlap window that looks like ``local gradients plus
the reducer bucket.''
\end{itemize}

\paragraph{ZeRO.}
ZeRO contributes three windows:
\begin{itemize}[leftmargin=1.5em]
\item a \textbf{fetch window}, which is the local trainable payload plus the
all-gather and prefetch buckets,
\item an \textbf{update window}, which is the optimizer update scratch,
\item a \textbf{communication window}, which is the local gradient payload plus
the reduce bucket.
\end{itemize}

The estimator then takes the largest of those windows instead of summing them.
That is the right abstraction for ``what is the largest temporary distributed
object live on this rank during the optimizer phase?''

The same idea in compact notation is:
\begin{align*}
W_{\mathrm{zero,fetch}} &=
M_{\mathrm{train}} + B_{\mathrm{zero,ag}} + B_{\mathrm{zero,pf}}, \\
W_{\mathrm{zero,comm}} &=
M_{\mathrm{grads}} + B_{\mathrm{zero,red}}, \\
W_{\mathrm{opt,zero}} &=
\max\!\left(
W_{\mathrm{zero,fetch}},
W_{\mathrm{update}},
B_{\mathrm{zero,red}}
\right).
\end{align*}

\subsection{Why One ZeRO-2 Full-FT Lifetime Path Remains}
One stage-specific path still remains for non-checkpointed \texttt{zero2
full\_ft}. It is not there because the code wants a special case for its own
sake. It is there because the canonical replay still shows a repeatable state
lifetime pattern in that regime.

The practical interpretation is:
\begin{itemize}[leftmargin=1.5em]
\item in backward, some gradient-related state behaves as if two copies of the
local gradient payload can overlap with the reduce bucket,
\item in optimizer step, the live window can include the master shard, local
gradients, update scratch, parameter partition support, and the ZeRO bucket
objects at the same time.
\end{itemize}

The report keeps this path explicit rather than pretending it comes from the
generic ZeRO window model. That is the more honest presentation.

The retained explicit lifetime path is:
\begin{align*}
W_{\mathrm{bwd}} &\leftarrow
\max\!\left(W_{\mathrm{bwd}}, 2 M_{\mathrm{grads}} + B_{\mathrm{zero,red}}\right), \\
\end{align*}
\begin{multline*}
W_{\mathrm{opt}} \leftarrow
\max\!\Big(
W_{\mathrm{opt}},
M_{\mathrm{master,shard}}
+ 2 M_{\mathrm{grads}}
+ W_{\mathrm{update}} \\
+ M_{\mathrm{param\_partition}}
+ B_{\mathrm{zero,ag}}
+ B_{\mathrm{zero,pf}}
+ B_{\mathrm{zero,red}}
\Big).
\end{multline*}

\subsection{Backward-End State}
The optimizer step does not start from a perfectly clean floor. Some state is
still live at the boundary between backward and the step. The estimator calls
that the \emph{backward-end state}.

The intuition is simple:
\begin{itemize}[leftmargin=1.5em]
\item single-GPU has no distributed bucket object to keep alive here,
\item DDP typically carries local gradients plus the reducer bucket,
\item ZeRO-2 carries local gradients plus the reduce bucket,
\item ZeRO-3 carries local gradients plus the all-gather-side support object
needed for the next step window.
\end{itemize}

This term matters because a lot of optimizer-step underestimation comes from
forgetting that the step does not begin from an empty temporary state.

The compact regime summary is:
\[
E_{\mathrm{bwd}} =
\begin{cases}
0 & \text{single\_gpu}, \\
M_{\mathrm{grads}} + B_{\mathrm{ddp}} & \text{ddp}, \\
M_{\mathrm{grads}} + B_{\mathrm{zero,red}} & \text{zero2}, \\
M_{\mathrm{grads}} + B_{\mathrm{zero,ag}} & \text{zero3}.
\end{cases}
\]

\section{Phase Assembly}
At this point the estimator has everything it needs:
\begin{itemize}[leftmargin=1.5em]
\item resident floors,
\item retained activations,
\item backward-end state,
\item phase-local workspace,
\item optimizer-step carryover.
\end{itemize}

The three candidate peaks are then:
\begin{align*}
M_{\mathrm{forward}} &= R_{\mathrm{base}} + A_{\mathrm{fwd}} + W_{\mathrm{fwd}}, \\
M_{\mathrm{backward}} &= R_{\mathrm{grad}} + A_{\mathrm{bwd}} + W_{\mathrm{bwd}}, \\
M_{\mathrm{optimizer}} &=
R_{\mathrm{grad}} + E_{\mathrm{bwd}} + W_{\mathrm{opt}} + C_{\mathrm{opt}}.
\end{align*}

\subsection{Optimizer Carryover}
Carryover is the least elegant part of the estimator, so it deserves a plain
language explanation.

\begin{quote}
Carryover means memory that is not part of the clean resident floor but still
has not been released when the optimizer step begins.
\end{quote}

Different regimes carry different things:
\begin{itemize}[leftmargin=1.5em]
\item some non-checkpointed ZeRO-2 full-FT paths carry a meaningful fraction of
the backward workspace or retained-forward state into the step,
\item checkpointed full-FT ZeRO paths mainly carry backward activation state,
\item checkpointed single-GPU LoRA can carry both retained-forward state and
some backward workspace into the step,
\item checkpointed distributed LoRA can also carry allocator reserve built up by
loss materialization and backward recompute, even when the clean resident and
activation terms look small at optimizer-step entry,
\item many simpler regimes carry essentially nothing extra.
\end{itemize}

The estimator keeps this logic as a short regime table in code because the
lifetimes really are regime-dependent. The important thing is that the report
states this openly instead of burying it in a single unexplained formula.

The newest example of that is checkpointed distributed LoRA. The current
heuristic adds the retained-forward proxy, one loss-workspace copy, the
backward workspace, and the small checkpoint-boundary or saved-input contexts
that tend to keep allocator reserve warm on long-context DDP and ZeRO runs.
That term is empirical rather than first-principles allocator simulation, but
it materially improves replay on the long-context OLMo 3 corpus without
changing the cleaner resident-state accounting elsewhere.

A compact generic view is:
\[
C_{\mathrm{opt}} = M_{\mathrm{carry,extra}},
\]
where $M_{\mathrm{carry,extra}}$ means ``state still live at optimizer-step
entry that has not already been counted in $R_{\mathrm{grad}} +
E_{\mathrm{bwd}}$.''

\subsection{Public Breakdown Semantics}
The public memory breakdown is aligned to the winning phase, not to every phase
at once.

That means:
\begin{itemize}[leftmargin=1.5em]
\item parameter, gradient, optimizer-state, and runtime-reserve terms always
show up directly,
\item activation and transient terms depend on which phase actually wins,
\item if optimizer step wins, the public activation bucket is intentionally set
to zero and the phase-local excess is shown as transient memory instead.
\end{itemize}

This choice keeps the public output readable while still allowing phase-aligned
calibration to compare like with like internally.

\section{Empirical Assumptions That Remain}
Most of the estimator is explicit tensor accounting. A few pieces still remain
as measured or runtime-driven assumptions:
\begin{itemize}[leftmargin=1.5em]
\item the attention workspace formulas are still proxies for kernel behavior,
not exact allocator traces,
\item runtime support is a named support model, not a full CUDA allocator
simulation,
\item checkpointed distributed LoRA optimizer carryover is still a calibrated
allocator-lifetime heuristic rather than a full reserve-pool simulator,
\item one stage-specific ZeRO-2 full-FT lifetime path remains because it still
improves replay quality materially.
\end{itemize}

Those assumptions are kept narrow on purpose. The estimator works best when the
core of the model stays about named objects and lifetimes, and only the truly
backend-shaped pieces remain empirical.
